\documentclass{article}
\usepackage{graphicx} % Required for inserting images
\usepackage{amsmath}
\usepackage{geometry}
\usepackage{hyperref}
\usepackage{graphicx}
\usepackage{parskip}

\geometry{margin=0.5in}

\title{Lagrangian decomposition with a streamfunction}
\author{jeffrey }
\date{October 2025}

\begin{document}

\maketitle

\section{Coupled asynchronous estimator with a spline-defined COM operator [Very good]}

We now derive the asynchronous analogue of the coupled system in Oscroft et al.
The center-of-mass (COM) trajectory is defined by a fast spline smoother, so the COM
evolution equation is obtained by applying that same smoothing operator to the full
fixed-frame dynamical model. Subtracting the COM equation from the full equation yields
the particle evolution in the COM frame.

The key point is that the mesoscale field remains present in \emph{both} equations.
The only mean-free assumption is that the submesoscale residual has zero projection onto
the COM spline space.

\subsection{Asynchronous observations}

Assume drifter $k$ is observed at times
\begin{equation}
t_{k,1},\dots,t_{k,N_k}, \qquad k=1,\dots,K,
\end{equation}
all lying in a common interval $[t_0,t_f]$, but not necessarily at simultaneous times.

Let
\begin{equation}
x_{k,n}=x_k(t_{k,n}), \qquad y_{k,n}=y_k(t_{k,n}),
\qquad
\dot{x}_{k,n}=\dot{x}_k(t_{k,n}), \qquad \dot{y}_{k,n}=\dot{y}_k(t_{k,n}),
\end{equation}
for $n=1,\dots,N_k$.

Let the total number of observations be
\begin{equation}
N_{\mathrm{tot}}=\sum_{k=1}^K N_k.
\end{equation}
Stack all observations in a single index $r=1,\dots,N_{\mathrm{tot}}$, with associated pair
$(k(r),n(r))$.

\subsection{Full fixed-frame model}

We write the fixed-frame velocity decomposition as
\begin{equation}
\dot{x}_k(t)=u_k^{\mathrm{meso}}(t)+u^{\mathrm{bg}}(t)+u_k^{\mathrm{sm}}(t),
\qquad
\dot{y}_k(t)=v_k^{\mathrm{meso}}(t)+v^{\mathrm{bg}}(t)+v_k^{\mathrm{sm}}(t).
\label{eq:full_fixed_model}
\end{equation}

Here
\begin{itemize}
    \item $(u_k^{\mathrm{meso}},v_k^{\mathrm{meso}})$ is the mesoscale velocity,
    \item $(u^{\mathrm{bg}},v^{\mathrm{bg}})$ is the background velocity, common to all drifters,
    \item $(u_k^{\mathrm{sm}},v_k^{\mathrm{sm}})$ is the submesoscale residual.
\end{itemize}

\subsection{Fast spline basis and COM smoothing operator}

Let the pooled observation times, sorted in nondecreasing order, be
\begin{equation}
\tau_1 \le \tau_2 \le \cdots \le \tau_{N_{\mathrm{tot}}}.
\end{equation}
Let $D(\tau_r)$ denote the number of deployed drifters at time $\tau_r$.
Define representative times recursively by
\begin{equation}
q_1=1,
\qquad
q_{m+1}=q_m+D(\tau_{q_m}),
\end{equation}
until $q_m>N_{\mathrm{tot}}$, and then set
\begin{equation}
\sigma_m=\tau_{q_m}, \qquad m=1,\dots,Q.
\end{equation}
These representative times define the knot sequence for a fast temporal spline basis
\begin{equation}
\{B_q(t)\}_{q=1}^{Q_b}.
\end{equation}
Here $Q_b$ is the number of fast temporal basis functions.


Define the fast-basis design matrix on the observation grid:
\begin{equation}
\mathcal{B}_{rq}=B_q\!\left(t_{k(r),n(r)}\right),
\qquad
r=1,\dots,N_{\mathrm{tot}}, \quad q=1,\dots,Q_b.
\label{eq:B_matrix}
\end{equation}

Stack the positions and velocities:
\begin{equation}
\mathbf{x}=
\begin{bmatrix}
x_{k(r),n(r)}
\end{bmatrix}_{r=1}^{N_{\mathrm{tot}}},
\qquad
\mathbf{y}=
\begin{bmatrix}
y_{k(r),n(r)}
\end{bmatrix}_{r=1}^{N_{\mathrm{tot}}},
\end{equation}
\begin{equation}
\dot{\mathbf{x}}=
\begin{bmatrix}
\dot{x}_{k(r),n(r)}
\end{bmatrix}_{r=1}^{N_{\mathrm{tot}}},
\qquad
\dot{\mathbf{y}}=
\begin{bmatrix}
\dot{y}_{k(r),n(r)}
\end{bmatrix}_{r=1}^{N_{\mathrm{tot}}}.
\end{equation}

The COM coefficients are obtained by ordinary least squares:
\begin{equation}
\hat{\mathbf{b}}^x=(\mathcal{B}^{\top}\mathcal{B})^{-1}\mathcal{B}^{\top}\mathbf{x},
\qquad
\hat{\mathbf{b}}^y=(\mathcal{B}^{\top}\mathcal{B})^{-1}\mathcal{B}^{\top}\mathbf{y}.
\end{equation}
Hence the COM trajectory on the observation grid is
\begin{equation}
\mathbf{m}_x=\mathcal{B}\hat{\mathbf{b}}^x,
\qquad
\mathbf{m}_y=\mathcal{B}\hat{\mathbf{b}}^y.
\end{equation}

Define the smoothing operator
\begin{equation}
S=\mathcal{B}(\mathcal{B}^{\top}\mathcal{B})^{-1}\mathcal{B}^{\top}.
\label{eq:S_operator}
\end{equation}
Let the derivative design matrix be
\begin{equation}
\dot{\mathcal{B}}_{rq}=B_q'\!\left(t_{k(r),n(r)}\right).
\end{equation}
Then
\begin{equation}
\mathbf{m}_x=S\mathbf{x},
\qquad
\mathbf{m}_y=S\mathbf{y},
\qquad
\dot{\mathbf{m}}_x=\dot{\mathcal{B}}\hat{\mathbf{b}}^x,
\qquad
\dot{\mathbf{m}}_y=\dot{\mathcal{B}}\hat{\mathbf{b}}^y.
\label{eq:S_com_relation}
\end{equation}
This is the quantity used in the current implementation: the COM velocity is obtained by
differentiating the fitted fast spline, while the projector $S$ is used for the COM-space
split of the mesoscale and residual velocities.

\subsection{Mesoscale streamfunction}

Let $\{T_\ell(t)\}_{\ell=1}^{L}$ be a slow temporal spline basis, and let
$X_i(\tilde{x})$, $Y_j(\tilde{y})$ be spatial basis functions in the COM frame.
We represent the mesoscale streamfunction as
\begin{equation}
\psi(\tilde{x},\tilde{y},t)
=
\sum_{i=1}^{n_x}\sum_{j=1}^{n_y}\sum_{\ell=1}^{L}
a_{ij\ell}\,
X_i(\tilde{x})\,Y_j(\tilde{y})\,T_\ell(t).
\label{eq:psi_model}
\end{equation}
The induced mesoscale velocity is
\begin{align}
u^{\mathrm{meso}}(\tilde{x},\tilde{y},t)
&=
-\sum_{i,j,\ell}
a_{ij\ell}\,
X_i(\tilde{x})\,Y_j'(\tilde{y})\,T_\ell(t),
\label{eq:u_meso_model}
\\
v^{\mathrm{meso}}(\tilde{x},\tilde{y},t)
&=
\sum_{i,j,\ell}
a_{ij\ell}\,
X_i'(\tilde{x})\,Y_j(\tilde{y})\,T_\ell(t).
\label{eq:v_meso_model}
\end{align}

Define the COM-frame coordinates at the observation times by
\begin{equation}
\tilde{x}_{k,n}=x_{k,n}-m_x(t_{k,n}),
\qquad
\tilde{y}_{k,n}=y_{k,n}-m_y(t_{k,n}).
\label{eq:tilde_coords}
\end{equation}

Now evaluate the mesoscale velocity at each observation and stack the results:
\begin{equation}
\mathbf{u}^{\mathrm{meso}}=
\begin{bmatrix}
u^{\mathrm{meso}}(\tilde{x}_{k(r),n(r)},\tilde{y}_{k(r),n(r)},t_{k(r),n(r)})
\end{bmatrix}_{r=1}^{N_{\mathrm{tot}}},
\end{equation}
\begin{equation}
\mathbf{v}^{\mathrm{meso}}=
\begin{bmatrix}
v^{\mathrm{meso}}(\tilde{x}_{k(r),n(r)},\tilde{y}_{k(r),n(r)},t_{k(r),n(r)})
\end{bmatrix}_{r=1}^{N_{\mathrm{tot}}}.
\end{equation}

Flattening the coefficients $a_{ij\ell}$ into a vector $\mathbf{a}$, we may write
\begin{equation}
\mathbf{u}^{\mathrm{meso}}=\mathbf{H}_x \mathbf{a},
\qquad
\mathbf{v}^{\mathrm{meso}}=\mathbf{H}_y \mathbf{a},
\label{eq:Hxy_def}
\end{equation}
where the rows of $\mathbf{H}_x$ and $\mathbf{H}_y$ are
\begin{equation}
(\mathbf{H}_x)_{r,ij\ell}
=
-\,X_i(\tilde{x}_{k(r),n(r)})\,Y_j'(\tilde{y}_{k(r),n(r)})\,T_\ell(t_{k(r),n(r)}),
\end{equation}
\begin{equation}
(\mathbf{H}_y)_{r,ij\ell}
=
\phantom{-}\,X_i'(\tilde{x}_{k(r),n(r)})\,Y_j(\tilde{y}_{k(r),n(r)})\,T_\ell(t_{k(r),n(r)}).
\end{equation}

\subsection{Background in the fast basis}

Represent the background velocity in the same fast basis:
\begin{equation}
u^{\mathrm{bg}}(t)=\sum_{q=1}^{Q_b} c_q^x B_q(t),
\qquad
v^{\mathrm{bg}}(t)=\sum_{q=1}^{Q_b} c_q^y B_q(t).
\label{eq:bg_model}
\end{equation}
Let
\begin{equation}
\mathbf{u}^{\mathrm{bg}}=\mathcal{B}\mathbf{c}^x,
\qquad
\mathbf{v}^{\mathrm{bg}}=\mathcal{B}\mathbf{c}^y.
\label{eq:bg_stacked}
\end{equation}

Because the background lies in the column space of $\mathcal{B}$,
\begin{equation}
S\mathbf{u}^{\mathrm{bg}}=\mathbf{u}^{\mathrm{bg}},
\qquad
S\mathbf{v}^{\mathrm{bg}}=\mathbf{v}^{\mathrm{bg}}.
\label{eq:bg_invariant}
\end{equation}
This follows because
\begin{equation}
S(\mathcal{B}\mathbf{c})
=
\mathcal{B}(\mathcal{B}^\top \mathcal{B})^{-1}\mathcal{B}^\top \mathcal{B}\mathbf{c}
=
\mathcal{B}\mathbf{c}
\end{equation}

\subsection{COM evolution equation}

Stacking \eqref{eq:full_fixed_model} gives
\begin{equation}
\dot{\mathbf{x}}
=
\mathbf{u}^{\mathrm{meso}}+\mathbf{u}^{\mathrm{bg}}+\mathbf{u}^{\mathrm{sm}},
\qquad
\dot{\mathbf{y}}
=
\mathbf{v}^{\mathrm{meso}}+\mathbf{v}^{\mathrm{bg}}+\mathbf{v}^{\mathrm{sm}}.
\label{eq:full_stacked}
\end{equation}

Applying the COM operator \(S\) and using \eqref{eq:S_com_relation}, we obtain
\begin{equation}
\dot{\mathbf{m}}_x
=
S\mathbf{H}_x\mathbf{a}
+
S\mathbf{u}^{\mathrm{bg}}
+
S\mathbf{u}^{\mathrm{sm}},
\qquad
\dot{\mathbf{m}}_y
=
S\mathbf{H}_y\mathbf{a}
+
S\mathbf{v}^{\mathrm{bg}}
+
S\mathbf{v}^{\mathrm{sm}}.
\end{equation}

Following the Oscroft logic, we impose that the submesoscale residual has no component
in the COM space,
\begin{equation}
S\mathbf{u}^{\mathrm{sm}}=\mathbf{0},
\qquad
S\mathbf{v}^{\mathrm{sm}}=\mathbf{0},
\label{eq:sm_zero_mean}
\end{equation}
so that the COM evolution reduces to
\begin{equation}
\dot{\mathbf{m}}_x
=
S\mathbf{H}_x\mathbf{a}
+
\mathbf{u}^{\mathrm{bg}},
\qquad
\dot{\mathbf{m}}_y
=
S\mathbf{H}_y\mathbf{a}
+
\mathbf{v}^{\mathrm{bg}},
\label{eq:com_equation_joint}
\end{equation}
where \(\mathbf{u}^{\mathrm{bg}}\) and \(\mathbf{v}^{\mathrm{bg}}\) are, by definition,
the residual background velocities in the COM space.

\subsection{Particle evolution in the COM frame}

Define the COM-frame velocities by
\begin{equation}
\dot{\tilde{\mathbf{x}}}=\dot{\mathbf{x}}-\dot{\mathbf{m}}_x,
\qquad
\dot{\tilde{\mathbf{y}}}=\dot{\mathbf{y}}-\dot{\mathbf{m}}_y.
\label{eq:tilde_vels}
\end{equation}

Subtracting \eqref{eq:com_equation_joint} from \eqref{eq:full_stacked} gives
\begin{equation}
\dot{\tilde{\mathbf{x}}}
=
(I-S)\mathbf{H}_x\mathbf{a}
+
\mathbf{u}^{\mathrm{sm}},
\qquad
\dot{\tilde{\mathbf{y}}}
=
(I-S)\mathbf{H}_y\mathbf{a}
+
\mathbf{v}^{\mathrm{sm}},
\label{eq:relative_equation_joint}
\end{equation}
where the submesoscale residual is mean-free by construction:
\begin{equation}
S\mathbf{u}^{\mathrm{sm}}=\mathbf{0},
\qquad
S\mathbf{v}^{\mathrm{sm}}=\mathbf{0}.
\end{equation}

\subsection{Mean-free assumption on the submesoscale residual}

We assume that the submesoscale residual has zero projection onto the
COM spline space:
\begin{equation}
S\mathbf{u}^{\mathrm{sm}}=\mathbf{0},
\qquad
S\mathbf{v}^{\mathrm{sm}}=\mathbf{0}.
\end{equation}
Thus the COM residual is identified with the background velocity, while
the relative residual is identified with the submesoscale velocity.
Equivalently,
\begin{equation}
\mathbf{u}^{\mathrm{bg}} = \dot{\mathbf{m}}_x - S\mathbf{H}_x\mathbf{a},
\qquad
\mathbf{v}^{\mathrm{bg}} = \dot{\mathbf{m}}_y - S\mathbf{H}_y\mathbf{a},
\label{eq:bg_residual_def}
\end{equation}
and
\begin{equation}
\mathbf{u}^{\mathrm{sm}} = \dot{\tilde{\mathbf{x}}} - (I-S)\mathbf{H}_x\mathbf{a},
\qquad
\mathbf{v}^{\mathrm{sm}} = \dot{\tilde{\mathbf{y}}} - (I-S)\mathbf{H}_y\mathbf{a}.
\label{eq:sm_residual_def}
\end{equation}

\subsection{Coupled least-squares system}

The unknown mesoscale coefficients \(\mathbf a\) are determined by minimizing
the residual energy in both the COM and relative equations.

Define
\begin{equation}
\mathbf{d}_{\mathrm{com}}
=
\begin{bmatrix}
\dot{\mathbf{m}}_x\\
\dot{\mathbf{m}}_y
\end{bmatrix},
\qquad
\mathbf{d}_{\mathrm{rel}}
=
\begin{bmatrix}
\dot{\tilde{\mathbf{x}}}\\
\dot{\tilde{\mathbf{y}}}
\end{bmatrix},
\end{equation}
and
\begin{equation}
\mathbf{M}_{\mathrm{com}}
=
\begin{bmatrix}
S\mathbf{H}_x\\
S\mathbf{H}_y
\end{bmatrix},
\qquad
\mathbf{M}_{\mathrm{rel}}
=
\begin{bmatrix}
(I-S)\mathbf{H}_x\\
(I-S)\mathbf{H}_y
\end{bmatrix}.
\end{equation}

Then
\begin{align}
\mathbf{d}_{\mathrm{com}}
&=
\mathbf{M}_{\mathrm{com}}\mathbf{a}
+
\boldsymbol{\epsilon}_{\mathrm{bg}},
\label{eq:com_ls_system}
\\
\mathbf{d}_{\mathrm{rel}}
&=
\mathbf{M}_{\mathrm{rel}}\mathbf{a}
+
\boldsymbol{\epsilon}_{\mathrm{sm}},
\label{eq:rel_ls_system}
\end{align}
where
\begin{equation}
\boldsymbol{\epsilon}_{\mathrm{bg}}
=
\begin{bmatrix}
\mathbf{u}^{\mathrm{bg}}\\
\mathbf{v}^{\mathrm{bg}}
\end{bmatrix},
\qquad
\boldsymbol{\epsilon}_{\mathrm{sm}}
=
\begin{bmatrix}
\mathbf{u}^{\mathrm{sm}}\\
\mathbf{v}^{\mathrm{sm}}
\end{bmatrix}.
\end{equation}

The mesoscale estimate is therefore obtained from
\begin{equation}
\hat{\mathbf a}
=
\arg\min_{\mathbf a}
\left(
\|
\mathbf{d}_{\mathrm{com}}-\mathbf{M}_{\mathrm{com}}\mathbf a
\|_2^2
+
\|
\mathbf{d}_{\mathrm{rel}}-\mathbf{M}_{\mathrm{rel}}\mathbf a
\|_2^2
\right).
\label{eq:joint_estimator}
\end{equation}

Equivalently, stacking the two systems gives
\begin{equation}
\underbrace{
\begin{bmatrix}
\mathbf{d}_{\mathrm{com}}\\
\mathbf{d}_{\mathrm{rel}}
\end{bmatrix}}_{\mathbf d}
=
\underbrace{
\begin{bmatrix}
\mathbf{M}_{\mathrm{com}}\\
\mathbf{M}_{\mathrm{rel}}
\end{bmatrix}}_{\mathbf X}
\mathbf a
+
\underbrace{
\begin{bmatrix}
\boldsymbol{\epsilon}_{\mathrm{bg}}\\
\boldsymbol{\epsilon}_{\mathrm{sm}}
\end{bmatrix}}_{\boldsymbol{\epsilon}}.
\label{eq:joint_stacked_system}
\end{equation}

\paragraph{Least-squares solution.}

The minimization problem \eqref{eq:joint_estimator} is a standard linear
least-squares problem with design matrix
\begin{equation}
\mathbf X =
\begin{bmatrix}
\mathbf{M}_{\mathrm{com}}\\
\mathbf{M}_{\mathrm{rel}}
\end{bmatrix},
\qquad
\mathbf d =
\begin{bmatrix}
\mathbf{d}_{\mathrm{com}}\\
\mathbf{d}_{\mathrm{rel}}
\end{bmatrix}.
\end{equation}
The normal equations are
\begin{equation}
\mathbf X^\top \mathbf X\, \hat{\mathbf a}
=
\mathbf X^\top \mathbf d.
\end{equation}
Using the block structure, this becomes
\begin{equation}
\left(
\mathbf{M}_{\mathrm{com}}^\top \mathbf{M}_{\mathrm{com}}
+
\mathbf{M}_{\mathrm{rel}}^\top \mathbf{M}_{\mathrm{rel}}
\right)\hat{\mathbf a}
=
\mathbf{M}_{\mathrm{com}}^\top \mathbf{d}_{\mathrm{com}}
+
\mathbf{M}_{\mathrm{rel}}^\top \mathbf{d}_{\mathrm{rel}}.
\label{eq:normal_equations}
\end{equation}

Substituting the definitions of the design matrices yields
\begin{align}
\mathbf{M}_{\mathrm{com}}^\top \mathbf{M}_{\mathrm{com}}
&=
\mathbf{H}_x^\top S \mathbf{H}_x
+
\mathbf{H}_y^\top S \mathbf{H}_y,
\\
\mathbf{M}_{\mathrm{rel}}^\top \mathbf{M}_{\mathrm{rel}}
&=
\mathbf{H}_x^\top (I-S) \mathbf{H}_x
+
\mathbf{H}_y^\top (I-S) \mathbf{H}_y.
\end{align}
Hence
\begin{equation}
\mathbf{X}^\top \mathbf{X}
=
\mathbf{H}_x^\top \mathbf{H}_x
+
\mathbf{H}_y^\top \mathbf{H}_y,
\end{equation}
because
\(
S + (I-S) = I.
\)

Similarly,
\begin{align}
\mathbf{X}^\top \mathbf d
&=
\mathbf{H}_x^\top S \dot{\mathbf m}_x
+
\mathbf{H}_y^\top S \dot{\mathbf m}_y
+
\mathbf{H}_x^\top (I-S)\dot{\tilde{\mathbf x}}
+
\mathbf{H}_y^\top (I-S)\dot{\tilde{\mathbf y}}.
\end{align}

Therefore the least-squares estimator is
\begin{equation}
\hat{\mathbf a}
=
\left(
\mathbf{H}_x^\top \mathbf{H}_x
+
\mathbf{H}_y^\top \mathbf{H}_y
\right)^{-1}
\left[
\mathbf{H}_x^\top S \dot{\mathbf m}_x
+
\mathbf{H}_y^\top S \dot{\mathbf m}_y
+
\mathbf{H}_x^\top (I-S)\dot{\tilde{\mathbf x}}
+
\mathbf{H}_y^\top (I-S)\dot{\tilde{\mathbf y}}
\right].
\label{eq:final_estimator}
\end{equation}

\subsection{Optional hard mesoscale constraints}

In the MATLAB implementation, the additive spatial gauge is first removed by
writing the full coefficient vector as
\begin{equation}
\mathbf c = G \boldsymbol{\alpha},
\end{equation}
where $G$ spans the gauge-free coefficient subspace. The unconstrained fit then
becomes the reduced least-squares problem
\begin{equation}
\hat{\boldsymbol{\alpha}}
=
\arg\min_{\boldsymbol{\alpha}}
\|
\mathbf d - \mathbf X G \boldsymbol{\alpha}
\|_2^2.
\end{equation}

The new optional constraint modes replace this with the equality-constrained
problem
\begin{equation}
\hat{\boldsymbol{\alpha}}
=
\arg\min_{\boldsymbol{\alpha}}
\|
\mathbf d - \mathbf X G \boldsymbol{\alpha}
\|_2^2
\qquad
\text{subject to}
\qquad
C\boldsymbol{\alpha}=0.
\label{eq:constrained_estimator}
\end{equation}
The constraint acts only on the mesoscale streamfunction coefficients; the
background and submesoscale components are still recovered afterward from the
COM and relative residuals.

For the \texttt{zeroVorticity} option, we enforce
\begin{equation}
\zeta = \psi_{qq} + \psi_{rr} = 0.
\end{equation}
For the \texttt{zeroStrain} option, we enforce
\begin{equation}
\psi_{qr} = 0,
\qquad
\psi_{qq} - \psi_{rr} = 0,
\end{equation}
which is equivalent to forcing both
\begin{equation}
\sigma_n = -2\psi_{qr} = 0,
\qquad
\sigma_s = \psi_{qq} - \psi_{rr} = 0.
\end{equation}

The implementation assembles these constraints by collocation over the
spline-support cells, not by adding a penalty term. For each nonzero knot span
in $q$, $r$, and $t$, it places $S_q+1$, $S_r+1$, and $S_t+1$ interior points,
forms the tensor-product support grid, and evaluates the derivative basis
matrices
\begin{equation}
A_{qq}, \qquad A_{rr}, \qquad A_{qr}
\end{equation}
on that grid. The reduced-space constraint matrix is then
\begin{equation}
C_{\zeta} = (A_{qq} + A_{rr})G
\end{equation}
for \texttt{zeroVorticity}, and
\begin{equation}
C_{\sigma}
=
\begin{bmatrix}
A_{qr}\\
A_{qq} - A_{rr}
\end{bmatrix}
G
\end{equation}
for \texttt{zeroStrain}. Redundant constraint rows are removed before solving
the equality-constrained normal equations.

Because each constrained diagnostic is piecewise polynomial on every
spline-support cell, vanishing on this collocation set forces the diagnostic to
vanish everywhere in the spline domain. In this sense the imposed constraint is
exact for the fitted spline model, not merely pointwise at the observations.

\paragraph{Interpretation.}
The \texttt{zeroVorticity} option forces the mesoscale streamfunction to be
harmonic. The \texttt{zeroStrain} option is much more restrictive: it removes
both strain components and leaves only the rigid-rotation-type spatial
mesoscale family.

\subsection{Recovered fields}

The fitted mesoscale velocity is
\begin{align}
\hat{u}^{\mathrm{meso}}(\tilde{x},\tilde{y},t)
&=
-\sum_{i,j,\ell}\hat{a}_{ij\ell}\,X_i(\tilde{x})Y_j'(\tilde{y})T_\ell(t),
\\
\hat{v}^{\mathrm{meso}}(\tilde{x},\tilde{y},t)
&=
\sum_{i,j,\ell}\hat{a}_{ij\ell}\,X_i'(\tilde{x})Y_j(\tilde{y})T_\ell(t).
\end{align}

The recovered background velocity is defined as the COM residual:
\begin{equation}
\hat{\mathbf{u}}^{\mathrm{bg}}
=
\dot{\mathbf{m}}_x
-
S\mathbf{H}_x\hat{\mathbf a},
\qquad
\hat{\mathbf{v}}^{\mathrm{bg}}
=
\dot{\mathbf{m}}_y
-
S\mathbf{H}_y\hat{\mathbf a}.
\end{equation}

The recovered submesoscale residual is defined as the relative residual:
\begin{equation}
\hat{\mathbf{u}}^{\mathrm{sm}}
=
\dot{\tilde{\mathbf{x}}}
-
(I-S)\mathbf{H}_x\hat{\mathbf a},
\qquad
\hat{\mathbf{v}}^{\mathrm{sm}}
=
\dot{\tilde{\mathbf{y}}}
-
(I-S)\mathbf{H}_y\hat{\mathbf a}.
\end{equation}
In the current implementation, $\hat{\mathbf{u}}^{\mathrm{bg}}$ and
$\hat{\mathbf{v}}^{\mathrm{bg}}$ are projected back into the fast spline basis before
integration, while the per-drifter mesoscale and submesoscale trajectory components are
rebuilt from interpolating splines on each drifter's native observation times.

\subsection{Trajectory decomposition in the fixed and COM frames}

The velocity decomposition \eqref{eq:full_fixed_model}, together with the
definitions of the recovered fields, induces a corresponding decomposition
of the trajectories by time integration. Because the observations are
asynchronous, the component trajectories must be defined with consistent
anchoring in time.

\paragraph{Background trajectory.}
The background velocity is defined as the residual of the COM equation,
\begin{equation}
\hat{\mathbf{u}}^{\mathrm{bg}}
=
\dot{\mathbf{m}}_x
-
S\mathbf{H}_x\hat{\mathbf a},
\qquad
\hat{\mathbf{v}}^{\mathrm{bg}}
=
\dot{\mathbf{m}}_y
-
S\mathbf{H}_y\hat{\mathbf a}.
\end{equation}
The corresponding background trajectory is defined as the anchored integral
\begin{equation}
x^{\mathrm{bg}}(t) := \int_{t_0}^{t} \hat{u}^{\mathrm{bg}}(s)\,ds,
\qquad
y^{\mathrm{bg}}(t) := \int_{t_0}^{t} \hat{v}^{\mathrm{bg}}(s)\,ds,
\end{equation}
with
\begin{equation}
x^{\mathrm{bg}}(t_0)=0,
\qquad
y^{\mathrm{bg}}(t_0)=0.
\end{equation}
In implementation, this residual is projected onto the fast spline basis
$\mathcal{B}$ and then integrated once to obtain a common smooth background path.
The fixed-frame background contribution for drifter $k$ is the reanchored path
\begin{equation}
x_k^{\mathrm{bg}}(t)=x^{\mathrm{bg}}(t)-x^{\mathrm{bg}}(t_{k,0}),
\qquad
y_k^{\mathrm{bg}}(t)=y^{\mathrm{bg}}(t)-y^{\mathrm{bg}}(t_{k,0}).
\end{equation}

\paragraph{Mesoscale trajectory.}
For each drifter $k$, the mesoscale trajectory contribution is defined by
\begin{equation}
x_k^{\mathrm{meso}}(t)
\;:=\;
x_k(t_{k,0})+\int_{t_{k,0}}^{t}
\hat{u}_k^{\mathrm{meso}}(s)\,ds,
\qquad
y_k^{\mathrm{meso}}(t)
\;:=\;
y_k(t_{k,0})+\int_{t_{k,0}}^{t}
\hat{v}_k^{\mathrm{meso}}(s)\,ds,
\end{equation}
where $t_{k,0}$ is the first observation time of drifter $k$, and
\begin{equation}
\hat{u}_k^{\mathrm{meso}}(t)
=
u^{\mathrm{meso}}(\tilde{x}_k(t),\tilde{y}_k(t),t),
\qquad
\hat{v}_k^{\mathrm{meso}}(t)
=
v^{\mathrm{meso}}(\tilde{x}_k(t),\tilde{y}_k(t),t).
\end{equation}
This definition ensures
\begin{equation}
x_k^{\mathrm{meso}}(t_{k,0})=x_k(t_{k,0}),
\qquad
y_k^{\mathrm{meso}}(t_{k,0})=y_k(t_{k,0}).
\end{equation}

\paragraph{Submesoscale trajectory.}
The submesoscale velocity is defined as the residual of the relative equation,
\begin{equation}
\hat{\mathbf{u}}^{\mathrm{sm}}
=
\dot{\tilde{\mathbf{x}}}
-
(I-S)\mathbf{H}_x\hat{\mathbf a},
\qquad
\hat{\mathbf{v}}^{\mathrm{sm}}
=
\dot{\tilde{\mathbf{y}}}
-
(I-S)\mathbf{H}_y\hat{\mathbf a}.
\end{equation}
For each drifter $k$, the corresponding trajectory contribution is
\begin{equation}
x_k^{\mathrm{sm}}(t)
:=
\int_{t_{k,0}}^{t} \hat{u}_k^{\mathrm{sm}}(s)\,ds,
\qquad
y_k^{\mathrm{sm}}(t)
:=
\int_{t_{k,0}}^{t} \hat{v}_k^{\mathrm{sm}}(s)\,ds,
\end{equation}
so that
\begin{equation}
x_k^{\mathrm{sm}}(t_{k,0})=0,
\qquad
y_k^{\mathrm{sm}}(t_{k,0})=0.
\end{equation}

\paragraph{Fixed-frame trajectory decomposition.}
With these definitions, the observed trajectory of drifter $k$ satisfies
\begin{align}
x_k(t)
&= 
x_k^{\mathrm{bg}}(t)
+
x_k^{\mathrm{meso}}(t)
+
x_k^{\mathrm{sm}}(t),
\\
y_k(t)
&=
y_k^{\mathrm{bg}}(t)
+
y_k^{\mathrm{meso}}(t)
+
y_k^{\mathrm{sm}}(t).
\end{align}
Differentiation recovers the velocity decomposition
\eqref{eq:full_fixed_model}.

\paragraph{COM-frame trajectory decomposition.}
Let $\tilde{x}_k = x_k - m_x$ and $\tilde{y}_k = y_k - m_y$ denote
centered coordinates. Then the background contribution cancels, and the
relative motion decomposes as
\begin{equation}
\dot{\tilde{\mathbf{x}}}
=
(I-S)\mathbf{H}_x \hat{\mathbf a}
+
\hat{\mathbf{u}}^{\mathrm{sm}},
\qquad
\dot{\tilde{\mathbf{y}}}
=
(I-S)\mathbf{H}_y \hat{\mathbf a}
+
\hat{\mathbf{v}}^{\mathrm{sm}}.
\end{equation}
Define the centered mesoscale contribution for drifter $k$ by
\begin{equation}
\tilde{x}_k^{\mathrm{meso}}(t)
=
\tilde{x}_k(t_{k,0})+\int_{t_{k,0}}^{t} \hat{\tilde{u}}_k^{\mathrm{meso}}(s)\,ds,
\qquad
\tilde{y}_k^{\mathrm{meso}}(t)
=
\tilde{y}_k(t_{k,0})+\int_{t_{k,0}}^{t} \hat{\tilde{v}}_k^{\mathrm{meso}}(s)\,ds,
\end{equation}
where
\begin{equation}
\hat{\tilde{\mathbf{u}}}^{\mathrm{meso}}=(I-S)\mathbf{H}_x\hat{\mathbf a},
\qquad
\hat{\tilde{\mathbf{v}}}^{\mathrm{meso}}=(I-S)\mathbf{H}_y\hat{\mathbf a}.
\end{equation}
The centered submesoscale contribution is zero-anchored:
\begin{equation}
\tilde{x}_k^{\mathrm{sm}}(t)
=
\int_{t_{k,0}}^{t} \hat{\tilde{u}}_k^{\mathrm{sm}}(s)\,ds,
\qquad
\tilde{y}_k^{\mathrm{sm}}(t)
=
\int_{t_{k,0}}^{t} \hat{\tilde{v}}_k^{\mathrm{sm}}(s)\,ds,
\end{equation}
we obtain
\begin{equation}
\tilde{x}_k(t)
=
\tilde{x}_k^{\mathrm{meso}}(t)
+
\tilde{x}_k^{\mathrm{sm}}(t),
\qquad
\tilde{y}_k(t)
=
\tilde{y}_k^{\mathrm{meso}}(t)
+
\tilde{y}_k^{\mathrm{sm}}(t).
\end{equation}
In the current code, these integrals are implemented drifter-by-drifter with
interpolating splines on the native sample times, using spline degree
$\min(3,N_k-1)$ for each velocity component.

\paragraph{Interpretation.}
In the fixed frame, the trajectory decomposes into background,
mesoscale, and submesoscale contributions. In the COM frame, the
background motion is removed, leaving a decomposition into mesoscale and
submesoscale relative motion. The background and submesoscale components
are defined as residuals in complementary subspaces, ensuring that the
mesoscale field captures as much variance as possible in both the COM
and relative dynamics.

\subsection{Zero gauge}

 The centered-frame streamfunction is represented as
 \[
 \psi(q,r,t) = \sum_{k=1}^{M_t} \sum_{j=1}^{M_r} \sum_{i=1}^{M_q}
 c_{ijk} B_i(q) B_j(r) T_k(t),
 \]
 where $B_i$, $B_j$, and $T_k$ are the tensor-product B-spline
 basis functions. Because the fit uses only $u = -\psi_r$ and
 $v = \psi_q$, the coefficients are unchanged if
 \[
 \psi(q,r,t) \mapsto \psi(q,r,t) + a(t).
 \]

 The implementation removes this gauge by constructing the spatial
 coefficient vector $\phi \in \mathbf{R}^{M_q M_r}$ that reproduces
 the constant field on a dense support grid,
\[
 A_{\mathrm{sp}} \phi = \mathbf{1},
 \]
 where $A_{\mathrm{sp}}$ is the spatial tensor-product basis evaluated
 on the $(q,r)$ support grid used to identify the constant mode.
 The retained spatial subspace is then
\[
 G_{\mathrm{sp}} = \operatorname{null}(\phi^\top),
 \]
 and the full reduced basis is
\[
 G = I_{M_t} \otimes G_{\mathrm{sp}},
 \qquad
 c = G \alpha.
 \]
 The fitted coefficients come from the reduced least-squares problem
\[
\hat{\alpha} = \arg\min_{\alpha} \|\mathbf{X} G \alpha - \mathbf{d}\|_2^2,
\qquad
\hat{c} = G \hat{\alpha},
\]
 with $\mathbf{X}$ and $\mathbf{d}$ defined in
\eqref{eq:joint_stacked_system}. This is exactly the gauge-reduced solve
 used in the current MATLAB implementation.

\end{document}
